\documentclass[11pt,a4paper]{article}

\usepackage[utf8]{inputenc}
\usepackage[T1]{fontenc}
\usepackage[ngerman]{babel}
\usepackage{amsmath,amssymb,amsthm}
\usepackage[a4paper,margin=2.5cm]{geometry}

\title{EABC-Operatoralgebra und Renormierungskette\\[0.3em]
\large Ikosaeder--Dodekaeder-Dualität}
\author{Thomas Hoffbauer}
\date{}

\newtheorem{theorem}{Satz}
\newtheorem{definition}{Definition}
\newtheorem{remark}{Bemerkung}
\newtheorem{proposition}{Proposition}

\newcommand{\R}{\mathbb{R}}
\newcommand{\Z}{\mathbb{Z}}

\begin{document}
\maketitle

%=============================================================================
\section{EABC-Familien und Anisotropietensor}
%=============================================================================

Die vier invertierbaren Restklassen modulo~$12$ definieren die EABC-Familien
\[
  E \equiv 1,\qquad A \equiv 5,\qquad B \equiv 7,\qquad C \equiv 11 \pmod{12},
\]
mit Gewichten $w_E=1$, $w_A=5$, $w_B=7$, $w_C=11$.

Sei $\{v_i\}_{i=1}^{12}\subset S^2\subset\R^3$ ein 12-Kuss-System (Einheitsrichtungen)
und $\sigma\colon\{1,\ldots,12\}\to\{E,A,B,C\}$ eine Label-Zuweisung.
Der \emph{Anisotropietensor} ist
\[
  M(\sigma)
  \;=\;
  \sum_{i=1}^{12} w_{\sigma(i)}\, v_i v_i^{\mathsf T}
  \;\in\;
  \mathrm{Sym}_3(\R).
\]

\begin{definition}[Ordnungsparameter]
\[
  \Delta(M) \;:=\; \lambda_{\max}(M) - \lambda_{\min}(M),
\]
wobei $\lambda_1\le\lambda_2\le\lambda_3$ die Eigenwerte von $M$ bezeichnen.
Isotropie entspricht $\Delta(M)=0$.
\end{definition}

%=============================================================================
\section{Renormierungsoperatoren}
%=============================================================================

Sei $K_n=(V_n,\sigma_n)$ eine Konfiguration aus Vertices $V_n$ und Zuweisung $\sigma_n$.

\begin{definition}[$R_{\mathrm{EABC}}$]
Der Operator $R_{\mathrm{EABC}}$ projiziert $\sigma$ auf die geometrisch zulässige
3+3+3+3-Zuweisung: pro Ecke $v_i$ wird die mod-$12$-Restklasse
$r(v_i)$ aus den skalierten Koordinaten und die nächste EABC-Restklasse
bestimmt. $R_{\mathrm{EABC}}$ ist idempotent; Fixpunkte erfordern perfekte
geometrische Konsistenz der Koordinaten-Klassifikation.
\end{definition}

\begin{definition}[$P_I$ (Voronoi-Pullback)]
$P_I$ bildet gestörte Kussrichtungen auf das Referenz-Ikosaeder ab:
\[
  P_I(v) \;=\; \arg\max_{r_j\in I_{12}} \langle v, r_j\rangle,
\]
wobei $I_{12}$ die 12 Referenz-Ecken des regulären Ikosaeders auf $S^2$ ist.
\end{definition}

\begin{definition}[$R^*_{\mathrm{EABC}}$]
Der robuste Fixpunktoperator ist die Komposition
\[
  R^*_{\mathrm{EABC}} \;:=\; P_I \circ R_{\mathrm{EABC}}.
\]
Zuerst geometrische Rückführung, dann Label-Projektion.
\end{definition}

\begin{proposition}[Idempotenz]
\[
  (R^*_{\mathrm{EABC}})^2 \;=\; R^*_{\mathrm{EABC}}.
\]
\end{proposition}

\begin{proof}[Skizze]
Nach $P_I$ liegen alle Vertices auf den Referenz-Ecken; $R_{\mathrm{EABC}}$ ist
auf dieser Untervarietät bereits Fixpunktabbildung. Ein zweiter Schritt ändert
weder Geometrie noch Labels.
\end{proof}

%=============================================================================
\section{Fixpunkt und Defektklassen}
%=============================================================================

\begin{definition}[Geometrischer Fixpunkt $I_{12}^{\mathrm{EABC}}$]
Die Konfiguration $I_{12}^{\mathrm{EABC}}$ besteht aus den 12 Ikosaeder-Ecken
mit der geometrischen EABC-Zuweisung (je 3 Labels pro Familie).
Am Fixpunkt gilt
\[
  \lambda_1 = \lambda_2 = \lambda_3 = 24, \qquad \Delta(M) = 0.
\]
\end{definition}

\begin{definition}[Defektklassen]
\begin{align*}
  \text{FIXPOINT}            &\colon \Delta(M)\approx 0,\; 3{+}3{+}3{+}3\text{-Balance},\; \text{ungestörte Geometrie},\\
  \text{LABEL\_DEFECT}       &\colon \text{Label-Zuweisung ohne 3+3+3+3-Balance (vor } R_{\mathrm{EABC}}\text{)},\\
  \text{DIRECTION\_DEFECT}   &\colon \text{balancierte Labels, aber }\Delta(M)>\varepsilon \text{ bei gestörter Geometrie},\\
  \text{CLASSIFICATION\_DEFECT} &\colon \text{mod-}12\text{-Balance der geometrischen Zuweisung gebrochen}.
\end{align*}
\end{definition}

%=============================================================================
\section{Dualität}
%=============================================================================

Ikosaeder--Dodekaeder-Dualität: $I_{12}\leftrightarrow D_{20}$ (12 Ecken $\leftrightarrow$ 20 Flächen).
Unter Dualitätsabbildung und Label-Pushforward gilt:

\begin{remark}[Dualitätsstabilität]
$\Delta(M)$ ist \emph{dualitätsstabil}: am geometrischen Fixpunkt stimmen
Ikosaeder- und Dodekaeder-Dual-Metrik überein.
Dagegen sind der Dipol $D=\sum_i w_{\sigma(i)} v_i$ und die Holonomie $H$
\emph{nicht} dualitätsinvariant.
\end{remark}

%=============================================================================
\section{Renormierungskette}
%=============================================================================

Die diskrete Kette lautet
\[
  K_n \;\longrightarrow\; K_{n+1}
  \;=\;
  \underbrace{\text{Volumen-Renormierung}}_{\lambda_{\mathrm{sphere}}}
  \;\circ\;
  \underbrace{R^*_{\mathrm{EABC}}}_{P_I \circ R_{\mathrm{EABC}}}
  \;\circ\;
  \underbrace{D_{20}}_{\text{Dual}}
  \;\circ\;
  \underbrace{I_{12}}_{\text{Ikosaeder}}
  \;\circ\;
  \underbrace{\text{12-Kuss}}_{\text{Einheitsrichtungen}}.
\]

\begin{definition}[Volumen-Renormierung]
Für äußeren Radius $R>0$ setze $R_{\mathrm{out}}=R(1+2\pi)$ und
\[
  \lambda_{\mathrm{sphere}}
  \;=\;
  \left(\frac{V_0}{V_{\mathrm{out}}}\right)^{\!1/3}
  \;=\;
  \frac{1}{1+2\pi},
\]
mit $V_0=\tfrac{4}{3}\pi R^3$, $V_{\mathrm{out}}=\tfrac{4}{3}\pi R_{\mathrm{out}}^3$.
Die Vertex-Skalierung erfolgt isotrop: $v_i\mapsto \lambda_{\mathrm{sphere}}\,v_i$.
\end{definition}

\begin{proposition}[Skalierungsinvarianz von $\Delta$ bei Isotropie]
Ist $\Delta(M)=0$, so bleibt $\Delta(M)$ unter isotroper Skalierung
$v_i\mapsto\lambda v_i$ invariant, da alle Eigenwerte mit $\lambda^2$ skaliert werden.
\end{proposition}

%=============================================================================
\section{Numerische Befunde}
%=============================================================================

Implementierung in \texttt{eabc\_icosahedron\_test.py}:

\begin{itemize}
\item $\varepsilon=0$: Nach $R_{\mathrm{EABC}}$ bzw.\ $R^*_{\mathrm{EABC}}$
  gilt $\Delta(M)\approx 0$ (Fixpunkt erreicht).
\item $\varepsilon>0$: Gestörte Geometrie ändert mod-$12$-Klassifikation;
  $R_{\mathrm{EABC}}$ allein reicht nicht --- gekoppelter Operator $R^*_{\mathrm{EABC}}$
  (Pullback + Label-Projektion) konvergiert zuverlässiger zu $\Delta(M)\approx 0$.
\item Gekoppelt vs.\ $R_{\mathrm{EABC}}$ allein: Bei $\varepsilon>0$ liefert
  $R^*_{\mathrm{EABC}}$ niedrigere mittlere $\Delta(M)$ und höhere Fixpunkt-Konvergenzrate.
\item Volumen-Schritt: $\lambda_{\mathrm{sphere}}=1/(1+2\pi)$; $\Delta(M)$ bleibt
  nach isotroper Skalierung null, wenn es vorher null war.
\end{itemize}

\end{document}
